% ─────────────────────────────────────────────────────────────────────────────
\section{Whitebox Robustness Under Direct Attacks}
\label{sec:ch6_whitebox}

\subsection{Aggregate Comparison Across Methods}
\label{subsec:ch6_aggregate_comparison}

\begin{table}[H]
\centering
\caption[Aggregate robust accuracy on the CIFAR-10 evaluation suite]{Aggregate robust accuracy (\%) on the CIFAR-10 evaluation suite,
reported as mean $\pm$ standard deviation across 20 random-seed runs.
Mean(8) is the average robust accuracy over the
eight non-AutoAttack evaluation attacks. Worst(8) is the lowest robust accuracy
among those same eight attacks. The
held-out attack mean averages FGSM-RS, TPGD, MI-FGSM, PGD-$\ell_2$,
DeepFool-$\ell_2$, and CW-$\ell_2$. The held-out $\ell_\infty$ mean averages
FGSM-RS, TPGD, and MI-FGSM; the held-out $\ell_2$ mean averages PGD-$\ell_2$,
DeepFool-$\ell_2$, and CW-$\ell_2$. The seen-source mean averages
PGD-$\ell_\infty$ and DDN-$\ell_2$. AutoAttack is reported separately and is
not included in Mean(8). Bold values indicate the best method for each metric.}
\label{tab:ch6_aggregate_results}
\footnotesize
\setlength{\tabcolsep}{1.25pt}
\renewcommand{\arraystretch}{1.08}
\begin{tabular}{@{}lcccc@{}}
\toprule
Metric & PGD-AT & DDN-AT & Multi-AT & AttackDRO++ (Ours) \\
\midrule
Mean(8)
& $60.499 \pm 0.327$
& $56.226 \pm 0.233$
& $62.045 \pm 0.447$
& \bestval{62.324 $\pm$ 0.147} \\

Held-out attack mean
& $62.268 \pm 0.390$
& $59.628 \pm 0.254$
& $64.083 \pm 0.505$
& \bestval{64.401 $\pm$ 0.155} \\

\hspace{1em}Held-out $\ell_\infty$ mean
& \bestval{62.515 $\pm$ 0.334}
& $51.594 \pm 0.306$
& $60.916 \pm 0.427$
& $61.149 \pm 0.170$ \\

\hspace{1em}Held-out $\ell_2$ mean
& $62.021 \pm 0.474$
& \bestval{67.661 $\pm$ 0.278}
& $67.249 \pm 0.603$
& $67.654 \pm 0.223$ \\

Seen-source mean
& $55.192 \pm 0.182$
& $46.020 \pm 0.247$
& $55.934 \pm 0.303$
& \bestval{56.093 $\pm$ 0.166} \\

\midrule
Worst(8)
& \bestval{49.656 $\pm$ 0.161}
& $25.957 \pm 0.378$
& $45.082 \pm 0.281$
& $45.079 \pm 0.237$ \\
\bottomrule



\end{tabular}
\end{table}
We compare four methods on 20 paired random-seed runs: PGD-AT,
DDN-AT, Multi-AT, and AttackDRO++ (Ours).
AttackDRO++ improves Mean(8) over Multi-AT by $+0.279$~pp ($p = 0.0076$,
14/20 wins), with the strongest channel on the held-out $\ell_2$ mean ($+0.405$~pp,
$p = 0.0079$, 16/20). It improves Mean(8) over PGD-AT by $+1.826$~pp and
over DDN-AT by $+6.098$~pp, both at $p < 10^{-10}$ on all 20 seeds, and
achieves $2.5$--$3.3{\times}$ tighter seed-to-seed standard deviation than Multi-AT on every average-case metric. Worst(8) and AutoAttack are
statistically indistinguishable from Multi-AT ($|\Delta| < 0.005$~pp,
$p > 0.97$): AttackDRO++ improves average and held-out robustness without
disturbing standardized worst-case behavior. The single-attack baselines
remain stronger than AttackDRO++ on the family they train against
(PGD-AT on the held-out $\ell_\infty$ mean, DDN-AT on the held-out $\ell_2$ mean) but lose
substantially on the other, confirming the specialist-versus-balanced
interpretation that motivates multi-attack training.

AttackDRO++ achieves the highest Mean(8), held-out attack mean, and
seen-source mean. PGD-AT wins the held-out $\ell_\infty$ mean and Worst(8);
DDN-AT wins the held-out $\ell_2$ mean. The bold pattern visualizes the specialist-versus-balanced
interpretation: each single-AT method dominates one column, while
AttackDRO++ is best on every column-aggregating metric.

\paragraph{Effect-size summary.}
Table~\ref{tab:ch6_effect_size_summary} reports paired effect magnitudes
between AttackDRO++ and each baseline. Cohen's $d_z = \Delta / \sigma_\Delta$
is the standardized paired effect size; conventional thresholds are
$0.2$ (small), $0.5$ (medium), $0.8$ (large). Statistical significance
is reported in Table~\ref{tab:ch6_paired_aggregate_all_baselines}.

\begin{table}[H]
\centering
\caption{Effect-size summary for AttackDRO++ (Ours) against each baseline across 20 paired random-seed runs.
$\Delta$ in percentage points; negative values mean the baseline is
stronger.}
\label{tab:ch6_effect_size_summary}
\small
\setlength{\tabcolsep}{4pt}
\renewcommand{\arraystretch}{1.10}
\begin{tabular}{lrrrrrr}
\toprule
& \multicolumn{2}{c}{vs.\ Multi-AT} & \multicolumn{2}{c}{vs.\ PGD-AT} & \multicolumn{2}{c}{vs.\ DDN-AT} \\
\cmidrule(lr){2-3}\cmidrule(lr){4-5}\cmidrule(lr){6-7}
Metric            & $\Delta$ (pp) & $d_z$    & $\Delta$ (pp) & $d_z$     & $\Delta$ (pp) & $d_z$ \\
\midrule
Mean(8)           & $+0.279$ & $0.67$   & $+1.826$ & $5.98$    & $+6.098$ & $23.09$ \\
Held-out attack mean & $+0.319$ & $0.67$   & $+2.134$ & $5.65$    & $+4.774$ & $16.91$ \\
\quad Held-out $\ell_\infty$ mean & $+0.233$ & $0.57$   & $-1.366$ & $-4.27$   & $+9.555$ & $25.44$ \\
\quad Held-out $\ell_2$ mean      & $+0.405$ & $0.66$   & $+5.633$ & $11.64$   & $-0.008$ & $-0.03$ \\
Seen-source mean  & $+0.159$ & $0.55$   & $+0.901$ & $4.59$    & $+10.073$ & $31.72$ \\
\midrule
Worst(8)          & $-0.003$ & $-0.01$  & $-4.576$ & $-15.36$  & $+19.122$ & $36.91$ \\
\bottomrule
\end{tabular}
\end{table}

Two regimes are visible. Against Multi-AT, every average-case metric is a
medium-effect gain ($d_z \in [0.55, 0.67]$, $\Delta \in [+0.16, +0.41]$~pp)
with Worst(8) at zero effect: AttackDRO++ shifts the average-case
distribution by a small but consistent amount without disturbing
worst-case behavior. Against the single-AT baselines, $d_z$ exceeds $4$
on most metrics but flips negative on the family each specialist trains
against. The medium-effect Multi-AT comparison validates the cluster-DRO
objective as a refinement of Multi-AT; the large-but-
asymmetric single-AT comparisons quantify the specialist tradeoff that
motivates the multi-attack setting in the first place.

\paragraph{Paired aggregate robustness comparisons.}
Table~\ref{tab:ch6_paired_aggregate_all_baselines} separates the main Multi-AT comparison from the
single-attack specialist comparisons. AutoAttack is excluded from this table and reported
separately in Section~\ref{subsec:ch6_hardest_evaluation_cases}.

\begin{table}[H]
\centering
\caption{Paired aggregate robustness comparison between AttackDRO++ (Ours) and each baseline across 20 paired random-seed runs.
\(\Delta\) is computed as AttackDRO++ minus the corresponding baseline in percentage
points. The \(p\)-value is from a two-sided paired \(t\)-test.}
\label{tab:ch6_paired_aggregate_all_baselines}
\footnotesize
\setlength{\tabcolsep}{3pt}
\renewcommand{\arraystretch}{1.08}
\begin{tabular}{llrrrrrr}
\toprule
Baseline & Metric & \(n\) & \(\Delta\) (pp) & \(\sigma_\Delta\) & Wins & \(d_z\) & \(p\) \\
\midrule

\multicolumn{8}{l}{\textit{AttackDRO++ (Ours) vs. Multi-AT}} \\
\midrule
 Multi-AT & Mean(8) & 20 & $+0.279$ & $0.417$ & $14/20$ & $0.668$ & $0.0076$ \\
 Multi-AT & Held-out attack mean & 20 & $+0.319$ & $0.478$ & $13/20$ & $0.667$ & $0.0076$ \\
 Multi-AT & Held-out $\ell_\infty$ mean & 20 & $+0.233$ & $0.412$ & $14/20$ & $0.566$ & $0.0204$ \\
 Multi-AT & Held-out $\ell_2$ mean & 20 & $+0.405$ & $0.609$ & $16/20$ & $0.664$ & $0.0079$ \\
 Multi-AT & Seen-source mean & 20 & $+0.159$ & $0.291$ & $15/20$ & $0.546$ & $0.0246$ \\
 Multi-AT & Worst(8) & 20 & $-0.003$ & $0.369$ & $9/20$ & $-0.008$ & $0.971$ \\

\midrule
\multicolumn{8}{l}{\textit{AttackDRO++ (Ours) vs. PGD-AT}} \\
\midrule
PGD-AT & Mean(8) & 20 & $+1.826$ & $0.305$ & $20/20$ & $5.980$ & $<10^{-10}$ \\
PGD-AT & Held-out attack mean & 20 & $+2.134$ & $0.378$ & $20/20$ & $5.650$ & $<10^{-10}$ \\
PGD-AT & Held-out $\ell_\infty$ mean & 20 & $-1.366$ & $0.320$ & $0/20$ & $-4.273$ & $<10^{-10}$ \\
PGD-AT & Held-out $\ell_2$ mean & 20 & $+5.633$ & $0.484$ & $20/20$ & $11.638$ & $<10^{-20}$ \\
PGD-AT & Seen-source mean & 20 & $+0.901$ & $0.197$ & $20/20$ & $4.585$ & $<10^{-10}$ \\
PGD-AT & Worst(8) & 20 & $-4.576$ & $0.298$ & $0/20$ & $-15.364$ & $<10^{-20}$ \\

\midrule
\multicolumn{8}{l}{\textit{AttackDRO++ (Ours) vs. DDN-AT}} \\
\midrule
DDN-AT & Mean(8) & 20 & $+6.098$ & $0.264$ & $20/20$ & $23.087$ & $<10^{-20}$ \\
DDN-AT & Held-out attack mean & 20 & $+4.774$ & $0.282$ & $20/20$ & $16.910$ & $<10^{-20}$ \\
DDN-AT & Held-out $\ell_\infty$ mean & 20 & $+9.555$ & $0.376$ & $20/20$ & $25.440$ & $<10^{-20}$ \\
DDN-AT & Held-out $\ell_2$ mean & 20 & $-0.008$ & $0.319$ & $10/20$ & $-0.025$ & $0.914$ \\
DDN-AT & Seen-source mean & 20 & $+10.073$ & $0.318$ & $20/20$ & $31.723$ & $<10^{-20}$ \\
DDN-AT & Worst(8) & 20 & $+19.122$ & $0.518$ & $20/20$ & $36.909$ & $<10^{-20}$ \\

\bottomrule
\end{tabular}
\end{table}



Against Multi-AT, AttackDRO++ gives
small but statistically supported gains on all average-case metrics, while Worst(8) is
unchanged. Against PGD-AT, AttackDRO++ improves Mean(8) and held-out $\ell_2$
robustness, but PGD-AT remains stronger on held-out $\ell_\infty$ and Worst(8).
Against DDN-AT, AttackDRO++ improves all $\ell_\infty$-related aggregates while matching
DDN-AT on held-out $\ell_2$. This again supports the main interpretation that
single-attack baselines are specialists, while AttackDRO++ provides better average
cross-attack balance.
% ─────────────────────────────────────────────────────────────────────────────
\begin{table}[H]
\centering
\caption{Multiple-comparison-corrected paired significance tests for AttackDRO++ (Ours)
against each baseline across $n=20$ paired seeds. $\Delta$ is computed as
AttackDRO++ minus the corresponding baseline in percentage points. The table
reports bootstrap 95\% confidence intervals, seed-level win counts, Cohen's paired
effect size $d_z$, paired $t$-test $p$-values, Wilcoxon signed-rank $p$-values,
and Holm-corrected Wilcoxon $p$-values. Holm correction is applied within each
baseline comparison across the seven reported metrics. AutoAttack is excluded
from this table and reported separately in Section~\ref{subsec:ch6_hardest_evaluation_cases}.}
\label{tab:ch6_corrected_paired_tests}
\scriptsize
\setlength{\tabcolsep}{2.5pt}
\renewcommand{\arraystretch}{1.08}
\resizebox{\textwidth}{!}{%
\begin{tabular}{llrrrrrrr}
\toprule
Baseline & Metric
& $\Delta$ (pp)
& 95\% CI (pp)
& Wins
& $d_z$
& $p_t$
& $p_W$
& $p_{W,\mathrm{Holm}}$ \\
\midrule

\multicolumn{9}{l}{\textit{AttackDRO++ (Ours) vs.\ Multi-AT}} \\
\midrule
 Multi-AT & Clean
& +0.539 & [+0.195, +0.903] & 13/20 & 0.642 & 0.0098 & 0.0441 & 0.1449 \\

 Multi-AT & Mean(8)
& +0.279 & [+0.108, +0.463] & 14/20 & 0.668 & 0.0076 & 0.0136 & 0.0953 \\

 Multi-AT & Held-out attack mean
& +0.319 & [+0.123, +0.527] & 13/20 & 0.667 & 0.0076 & 0.0192 & 0.0962 \\

 Multi-AT & Held-out $\ell_\infty$ mean
& +0.233 & [+0.060, +0.410] & 14/20 & 0.566 & 0.0204 & 0.0362 & 0.1449 \\

 Multi-AT & Held-out $\ell_2$ mean
& +0.405 & [+0.144, +0.661] & 16/20 & 0.664 & 0.0079 & 0.0136 & 0.0953 \\

 Multi-AT & Seen-source mean
& +0.159 & [+0.038, +0.285] & 15/20 & 0.546 & 0.0246 & 0.0419 & 0.1449 \\

 Multi-AT & Worst(8)
& -0.003 & [-0.168, +0.155] & 9/20 & -0.008 & 0.9714 & 0.9839 & 0.9839 \\

\midrule
\multicolumn{9}{l}{\textit{AttackDRO++ (Ours) vs.\ PGD-AT}} \\
\midrule
PGD-AT & Clean
& +4.015 & [+3.666, +4.386] & 20/20 & 4.779 & $<0.0001$ & $<0.0001$ & $<0.0001$ \\

PGD-AT & Mean(8)
& +1.826 & [+1.702, +1.965] & 20/20 & 5.980 & $<0.0001$ & $<0.0001$ & $<0.0001$ \\

PGD-AT & Held-out attack mean
& +2.134 & [+1.980, +2.303] & 20/20 & 5.650 & $<0.0001$ & $<0.0001$ & $<0.0001$ \\

PGD-AT & Held-out $\ell_\infty$ mean
& -1.366 & [-1.496, -1.223] & 0/20 & -4.273 & $<0.0001$ & $<0.0001$ & $<0.0001$ \\

PGD-AT & Held-out $\ell_2$ mean
& +5.633 & [+5.441, +5.849] & 20/20 & 11.638 & $<0.0001$ & $<0.0001$ & $<0.0001$ \\

PGD-AT & Seen-source mean
& +0.901 & [+0.815, +0.984] & 20/20 & 4.585 & $<0.0001$ & $<0.0001$ & 0.0002 \\

PGD-AT & Worst(8)
& -4.576 & [-4.708, -4.454] & 0/20 & -15.364 & $<0.0001$ & $<0.0001$ & 0.0002 \\

\midrule
\multicolumn{9}{l}{\textit{AttackDRO++ (Ours) vs.\ DDN-AT}} \\
\midrule
DDN-AT & Clean
& -4.563 & [-4.712, -4.414] & 0/20 & -13.134 & $<0.0001$ & $<0.0001$ & $<0.0001$ \\

DDN-AT & Mean(8)
& +6.098 & [+5.986, +6.209] & 20/20 & 23.087 & $<0.0001$ & $<0.0001$ & $<0.0001$ \\

DDN-AT & Held-out attack mean
& +4.774 & [+4.654, +4.893] & 20/20 & 16.910 & $<0.0001$ & $<0.0001$ & $<0.0001$ \\

DDN-AT & Held-out $\ell_\infty$ mean
& +9.555 & [+9.399, +9.718] & 20/20 & 25.440 & $<0.0001$ & $<0.0001$ & $<0.0001$ \\

DDN-AT & Held-out $\ell_2$ mean
& -0.008 & [-0.143, +0.127] & 10/20 & -0.025 & 0.9137 & 0.9563 & 0.9563 \\

DDN-AT & Seen-source mean
& +10.072 & [+9.938, +10.208] & 20/20 & 31.723 & $<0.0001$ & $<0.0001$ & 0.0002 \\

DDN-AT & Worst(8)
& +19.122 & [+18.904, +19.344] & 20/20 & 36.909 & $<0.0001$ & $<0.0001$ & $<0.0001$ \\

\bottomrule
\end{tabular}%
}
\end{table}

Table~\ref{tab:ch6_paired_aggregate_all_baselines} reports the paired $t$-test
comparison. To align with the statistical protocol described in Chapter~5,
Table~\ref{tab:ch6_corrected_paired_tests} additionally reports Wilcoxon signed-rank
tests with Holm correction, bootstrap confidence intervals, paired effect sizes,
and seed-level win counts.
The multiple-comparison-corrected tests provide a conservative check on the
paired $t$-test results. Against Multi-AT, the key question is whether the
small positive gains on average-case metrics remain supported after Wilcoxon
testing and Holm correction. Clean accuracy is included as a sanity metric, but
the main robustness claims are based on Mean(8), held-out metrics, seen-source mean, and
Worst(8). Against the single-attack baselines, the same analysis mainly verifies
the specialist-versus-balanced pattern: PGD-AT remains stronger on
$\ell_\infty$-dominated metrics, DDN-AT remains competitive on held-out $\ell_2$,
and AttackDRO++ provides stronger average cross-attack balance.

After Holm correction, the Multi-AT comparison should be interpreted cautiously.
AttackDRO++ shows small positive differences over Multi-AT on Clean, Mean(8),
held-out attack mean, held-out $\ell_\infty$ mean, held-out $\ell_2$ mean, and
seen-source mean, and these
differences are supported by the paired $t$-tests, raw Wilcoxon tests, confidence
intervals, and moderate paired effect sizes. However, the Holm-corrected
Wilcoxon $p$-values for these metrics do not remain below the $0.05$ threshold.
Therefore, relative to Multi-AT, the evidence is best described as a small
positive trend rather than a uniformly correction-robust improvement. Worst(8)
is essentially unchanged, indicating that AttackDRO++ preserves the worst-case
attack performance of Multi-AT rather than improving it.

For the single-attack baselines, the corrected tests support the expected
specialist-versus-balanced interpretation. Compared with PGD-AT, AttackDRO++
achieves correction-robust gains on Mean(8), held-out attack mean, held-out $\ell_2$ mean, and
seen-source mean, but remains weaker on the held-out $\ell_\infty$ mean and Worst(8), consistent
with PGD-AT's specialization toward $\ell_\infty$ robustness. Compared with
DDN-AT, AttackDRO++ obtains correction-robust gains on Mean(8), held-out attack mean,
held-out $\ell_\infty$ mean, seen-source mean, and Worst(8), while the held-out $\ell_2$ mean is
statistically indistinguishable. Overall, Table~\ref{tab:ch6_corrected_paired_tests}
suggests that AttackDRO++ mainly improves average cross-attack balance and
held-out robustness, while its worst-case behavior depends on the baseline being
compared.


\subsection{Per-Attack Accuracy: Where Do Methods Diverge?}
\label{subsec:ch6_per_attack_accuracy}

The aggregate metrics show the overall trend, but they do not reveal which attacks
drive the improvement. Table~\ref{tab:ch6_per_attack_all_baselines} reports per-attack
paired comparisons between AttackDRO++ and each baseline. The table is divided into
three panels: the main comparison against Multi-AT, followed by the two
single-attack specialist comparisons.

\begin{table}[H]
\centering
\caption{Per-attack robust accuracy (\%) for all four methods across 20 paired random-seed runs.
Bold values indicate the best method for each attack. The right-hand
$\Delta$ and $p$ columns report the paired comparison of AttackDRO++ (Ours) against
Multi-AT. AttackDRO++ (Ours) wins the highest
score on 4 of 8 attacks; the two single-attack baselines together account
for the remaining 4, each winning on the family they were trained against.
$\Delta$ values use the bold-marker convention $^*$ for $p<0.05$ and
$^{**}$ for $p<0.01$.}
\label{tab:ch6_per_attack_all_baselines}
\footnotesize
\setlength{\tabcolsep}{3pt}
\renewcommand{\arraystretch}{1.10}
\resizebox{\textwidth}{!}{%
\begin{tabular}{llcccc|rr}
\toprule
Attack & Family & PGD-AT & DDN-AT & Multi-AT & AttackDRO++ (Ours) & $\Delta$ vs.\ Multi-AT & $p$ \\
\midrule
FGSM-RS         & held-out $\ell_\infty$
                & $67.896$
                & $64.834$
                & $68.483$
                & \bestval{68.939}
                & $+0.456^{**}$
                & $0.0034$ \\

PGD-$\ell_\infty$ & seen $\ell_\infty$
                & \bestval{49.656}
                & $25.957$
                & $45.082$
                & $45.079$
                & $-0.003$
                & $0.971$ \\

TPGD            & held-out $\ell_\infty$
                & \bestval{68.388}
                & $59.310$
                & $66.934$
                & $67.043$
                & $+0.109$
                & $0.513$ \\

MI-FGSM         & held-out $\ell_\infty$
                & \bestval{51.260}
                & $30.638$
                & $47.331$
                & $47.466$
                & $+0.135$
                & $0.116$ \\

\midrule
PGD-$\ell_2$    & held-out $\ell_2$
                & $64.319$
                & \bestval{69.581}
                & $69.197$
                & $69.571$
                & $+0.374^{*}$
                & $0.0120$ \\

DDN-$\ell_2$    & seen $\ell_2$
                & $60.728$
                & $66.083$
                & $66.785$
                & \bestval{67.107}
                & $+0.321^{*}$
                & $0.0136$ \\

DeepFool-$\ell_2$ & held-out $\ell_2$
                & $62.533$
                & \bestval{68.290}
                & $67.301$
                & $67.805$
                & $+0.504^{**}$
                & $0.0082$ \\

CW-$\ell_2$     & held-out $\ell_2$
                & $59.211$
                & $65.113$
                & $65.249$
                & \bestval{65.585}
                & $+0.336^{**}$
                & $0.0097$ \\
\bottomrule
\end{tabular}%
}
\end{table}
Table~\ref{tab:ch6_per_attack_all_baselines} disaggregates the Mean(8) result
into the eight evaluation attacks. AttackDRO++ achieves the highest
per-attack accuracy on four of eight attacks, including all three held-out
$\ell_2$ attacks where it beats Multi-AT at $p<0.05$. PGD-AT wins on
the three held-out $\ell_\infty$ attacks and on the seen-source attack
PGD-$\ell_\infty$ that it trains directly against; DDN-AT wins on
PGD-$\ell_2$ and DeepFool-$\ell_2$, the held-out $\ell_2$ attacks closest to
its training distribution. The pattern is consistent with a specialist
interpretation of single-attack training: each single-AT method dominates a
narrow attack family at the cost of substantial regression on the other.
AttackDRO++'s improvement over Multi-AT is largest on held-out
$\ell_2$ attacks ($+0.374$ to $+0.504$~pp, all $p<0.05$) and smallest on the
seen $\ell_\infty$ attack ($-0.003$~pp), which explains why Worst(8) is
preserved rather than improved.


\subsection{Robustness Under the Hardest Evaluation Cases}
\label{subsec:ch6_hardest_evaluation_cases}

This subsection examines two complementary stress indicators: the lowest robust accuracy across the eight main evaluation attacks and the separate AutoAttack-$\ell_\infty$ evaluation.

Two metrics test standardized worst-case behavior. Worst(8) is the minimum
robust accuracy across the eight evaluation attacks per checkpoint;
AutoAttack-$\ell_\infty$ is an ensemble $\ell_\infty$ stress test following
Croce and Hein \cite{croce2020reliable}. Tables~\ref{tab:ch6_worst_autoattack_sanity} and
\ref{tab:ch6_autoattack_fulltest} report each metric separately because they have
different evaluation panels: Worst(8) is computed across the full $n=20$
seed panel, while AutoAttack is reported in two layers --- a 512-sample
sanity evaluation across all $n=20$ seeds and a full-test evaluation on
the selected $n=5$ seed panel.

\begin{table}[H]
\centering
\caption{Worst(8) and 512-sample AutoAttack sanity results across
20 paired random-seed runs. Worst(8) is the minimum robust accuracy over the eight
non-AutoAttack evaluation attacks. For AttackDRO++ $-$ Multi-AT,
the paired differences are $-0.003$~pp on Worst(8) ($p=0.971$) and
$-0.001$~pp on AutoAttack ($p=0.998$). Bold values indicate the best method
for each metric.}
\label{tab:ch6_worst_autoattack_sanity}
\small
\setlength{\tabcolsep}{8pt}
\renewcommand{\arraystretch}{1.08}
\begin{tabular}{lrrrr}
\toprule
& \multicolumn{2}{c}{Worst(8)} & \multicolumn{2}{c}{AutoAttack-$\ell_\infty$} \\
\cmidrule(lr){2-3}\cmidrule(lr){4-5}
Method & Mean (\%) & $\sigma$ (pp) & Mean (\%) & $\sigma$ (pp) \\
\midrule
PGD-AT & \bestval{49.656} & $0.161$ & \bestval{44.912} & $1.136$ \\
DDN-AT      & $25.957$ & $0.378$ & $23.809$ & $0.796$ \\
Multi-AT                    & $45.082$ & $0.281$ & $41.318$ & $1.167$ \\
AttackDRO++ (Ours)          & $45.079$ & $0.237$ & $41.317$ & $1.050$ \\
\bottomrule
\end{tabular}
\end{table}


\begin{table}[H]
\centering
\caption{Full-test AutoAttack-$\ell_\infty$ results on the selected
$n=5$ paired seed panel. The paired difference is
AttackDRO++ $-$ Multi-AT: $-0.006$~pp, with $p=0.9584$.}
\label{tab:ch6_autoattack_fulltest}
\small
\setlength{\tabcolsep}{10pt}
\renewcommand{\arraystretch}{1.08}
\begin{tabular}{lrr}
\toprule
Method & Mean (\%) & $\sigma$ (pp) \\
\midrule
Multi-AT & $42.47$ & $0.30$ \\
AttackDRO++ (Ours) & $42.46$ & $0.20$ \\
\bottomrule
\end{tabular}
\end{table}

\paragraph{AttackDRO++ vs.\ Multi-AT.}
AttackDRO++ is statistically indistinguishable from Multi-AT on both
worst-case metrics. On the $n=20$ sanity panel, the paired differences are
$-0.003$~pp for Worst(8) ($p=0.971$) and $-0.001$~pp for AutoAttack
($p=0.998$). The full-test AutoAttack panel gives the same conclusion
($-0.006$~pp, $p=0.9584$). Thus, the claim is preservation rather than
improvement: AttackDRO++ improves average and held-out robustness without
degrading standardized worst-case behavior.

\paragraph{PGD-AT specialization.}
PGD-AT obtains the highest Worst(8) and AutoAttack because both metrics are
dominated by $\ell_\infty$ robustness, the threat family it trains against.
This advantage is therefore structural rather than evidence of broader
robustness. As Table~\ref{tab:ch6_per_attack_all_baselines} shows, PGD-AT loses
$5$--$6$~pp to AttackDRO++ on held-out $\ell_2$ attacks, while DDN-AT shows the
opposite failure mode by collapsing under $\ell_\infty$ stress. Single-attack
training therefore produces specialists; AttackDRO++ does not match PGD-AT's
$\ell_\infty$ depth, but provides the best average cross-attack balance.

\subsection{Training Stability: How Predictable Is the Outcome?}
\label{subsec:ch6_training_stability}

In addition to improving the mean of average-case robustness metrics, AttackDRO++ also
reduces seed-to-seed variation relative to Multi-AT. Table~\ref{tab:ch6_variance_reduction}
compares the standard deviation of each seed-level metric between the two methods.

\begin{table}[H]
\centering
\caption{Variance reduction across 20 paired random-seed runs. \(\sigma\) is the
standard deviation of the per-seed metric. AttackDRO++ substantially reduces
seed-to-seed variation on Mean(8), held-out averages, and held-out norm-specific
metrics, while Worst(8) and AutoAttack standard deviations remain similar across
methods.}
\label{tab:ch6_variance_reduction}
\small
\setlength{\tabcolsep}{5pt}
\renewcommand{\arraystretch}{1.08}
\begin{tabular}{lrrr}
\toprule
Metric & \(\sigma\) Multi-AT (pp) & \(\sigma\) AttackDRO++ (pp) & Ratio \\
\midrule
Mean(8)
& $0.447$ & $0.148$ & $3.03\times$ \\

Held-out attack mean
& $0.505$ & $0.155$ & $3.26\times$ \\

Held-out $\ell_2$ mean
& $0.603$ & $0.223$ & $2.70\times$ \\

Held-out $\ell_\infty$ mean
& $0.427$ & $0.170$ & $2.51\times$ \\

Seen-source mean
& $0.303$ & $0.166$ & $1.83\times$ \\

\midrule
Worst(8)
& $0.281$ & $0.237$ & $1.18\times$ \\

AutoAttack-$\ell_\infty$
& $1.167$ & $1.050$ & $1.11\times$ \\
\bottomrule
\end{tabular}
\end{table}

The reduction is strongest for held-out robustness. The standard deviation decreases by
\(3.26\times\) for the held-out attack mean, \(2.70\times\) for the held-out \(\ell_2\) mean, and
\(2.51\times\) for the held-out \(\ell_\infty\) mean. Mean(8) also becomes substantially more
stable, decreasing from \(0.447\)~pp under Multi-AT to \(0.148\)~pp under
AttackDRO++, a \(3.03\times\) reduction.

This suggests that the anchored Cluster-DRO objective not only improves the mean of
average-case robustness metrics, but also makes the Multi-AT training outcome less
sensitive to random seed. The effect is weaker for Worst(8) and AutoAttack, whose
standard deviations remain similar across methods. This is consistent with the previous
sections: AttackDRO++ preserves the hardest \(\ell_\infty\) stress metrics relative to
 Multi-AT, while its main benefit is on average and held-out robustness.

The whitebox results give the main aggregate evidence: AttackDRO++ improves
average and held-out robustness relative to Multi-AT while preserving the hardest
standardized checks. However, aggregate and per-attack metrics can still hide
which CIFAR-10 classes and attack--class cells drive the remaining failures. The
next section therefore drills down into class-wise whitebox robustness to test
whether the observed gains are broad or concentrated in a small part of the
evaluation grid.

